
Reinforcement Learning from Human Feedback assumes that human annotators provide a stable preference signal, yet the annotators themselves may adapt toward the AI style they repeatedly evaluate. This paper investigates this \textit{annotation drift} phenomenon in two RLHF preference datasets — Anthropic HH-RLHF (160,800 comparisons) and OpenAI WebGPT (19,578 comparisons) — using the Alignment Asymmetry Index (AAI), a composite instrument for measuring directional stylistic adaptation in preference labels. Components 1 and 2 of the AAI are evaluated in this work; the reward-model behavioral divergence component is fully specified but was not executed.

The primary finding is that verbosity preference reverses sign across annotation strata in HH-RLHF: early-round annotators are associated with a negative verbosity coefficient (β_L = −0.025, 95\% CI [−0.043, −0.006]), while later-round annotators are associated with a positive verbosity coefficient (β_L = +0.056, 95\% CI [+0.043, +0.068]), yielding Δβ_L = +0.080 with non-overlapping bootstrap confidence intervals. A secondary analysis on WebGPT using a between-group tercile design — adopted as a fallback due to the absence of worker identity metadata in the public dataset release — yields a statistically significant tercile F-statistic (F = 33.226, p ≈ 8.3×10⁻⁹) and a discriminant validity result confirming AI-typicality specificity (placebo empirical p = 0.48). Pre-registered effect-size thresholds were not met in the WebGPT analysis, and the planned within-annotator panel regression could not be executed.

Null results — including an untestable ambiguity-modulation interaction (p = 1.0, due to absent per-prompt disagreement labels) and subthreshold effects for hedging and structured reasoning — characterize the minimum data requirements for confirming individual annotator causal adaptation. These findings suggest that multi-round RLHF annotation datasets may encode temporal non-stationarity in the preference signal, motivating drift-aware quality auditing prior to reward model training.

